\documentclass{article}
\usepackage{PRIMEarxiv} % Core package for the PrimeAI Template
\usepackage{amsmath, amssymb, amsthm, bm, dcolumn} % Add amsthm here for the proof environment
\usepackage[numbers,sort&compress]{natbib} % Natbib for citations
\usepackage{graphicx} % For high-quality images
\usepackage{multicol} % Multi-column figures/tables
\usepackage{hyperref} % Hyperlinks
\usepackage{listings} % Code listings
\usepackage{authblk} % For structured affiliations
% Define theorem style (optional)
\theoremstyle{plain}
\newtheorem{theorem}{Theorem}
\renewcommand{\qedsymbol}{}

% Custom commands for primes and qubit encoding
\newcommand{\primeQubit}[1]{|p_{#1}\rangle}
\newcommand{\primeGate}[1]{U_{p_{#1}}}

\usepackage{wrapfig}
\usepackage[pscoord]{eso-pic}
\usepackage[fulladjust]{marginnote}
\reversemarginpar

% Typesetting improvements without footnote patching
\usepackage[protrusion=true, expansion=true, tracking=false]{microtype}
\microtypecontext{spacing=nonfrench}

% Line numbers
\usepackage[right]{lineno}

% Text layout - adjust as needed
\raggedright
\setlength{\parindent}{0.5cm}
\textwidth 5.25in 
\textheight 8.75in

% Set double spacing
\usepackage{setspace} 
\doublespacing

% Adjust width for specific content
\usepackage{changepage}

% Adjust caption style
\usepackage[aboveskip=1pt,labelfont=bf,labelsep=period,singlelinecheck=off]{caption}

% Remove brackets from references
\makeatletter
\renewcommand{\@biblabel}[1]{\quad#1.}
\makeatother

% Header, footer, and page numbers
\usepackage{lastpage,fancyhdr}
\pagestyle{fancy}
\fancyhf{}
\fancyhead[L]{Preprint - PrimeAI Enhanced Template}
\fancyfoot[C]{\scriptsize Multiplicity Theory © 2025 Ryan O. Van Gelder - Citizen Gardens \\ Licensed Under MIT and CC BY-NC-SA 4.0.}
\fancyfoot[R]{Page \thepage\ of \pageref{LastPage}}
\renewcommand{\footrule}{\hrule height 2pt \vspace{2mm}}

\begin{document}
\title{Quantum Anthropology and Social Physics through Multiplicity Theory}

\author{Ryan O. Van Gelder}
\affil{Citizen Gardens - The Foundation of Multiplicity}
\date{\today}
\maketitle

\begin{abstract}
This article explores the novel integration of Multiplicity Theory into the domains of Quantum Anthropology and Social Physics. It examines whether social structures and cultural phenomena can be modeled as quantum-like systems, leveraging principles such as entanglement, superposition, and tunneling. The paper develops mathematical frameworks to analyze how minor behavioral changes propagate through networks to drive macro-level societal shifts, incorporating experimental applications in real-world contexts.
\end{abstract}

\keywords{Multiplicity Theory, Quantum Anthropology, Social Physics, Entanglement, Superposition, Tunneling, Ethical Frameworks, Tensor Networks, Recursive Feedback, Stochastic Dynamics, Eigenvalue Analysis, Prime-Based Encoding, Societal Dynamics, Computational Sociology}

\section{Introduction}
Multiplicity Theory provides a robust mathematical foundation for understanding interconnected systems, blending classical and quantum principles. By examining human networks through this lens, we aim to:
\begin{itemize}
    \item Model shared decisions (entanglement), undecided outcomes (superposition), and unexpected transitions (tunneling).
    \item Quantify the propagation of micro-level changes through social networks.
    \item Design ethical frameworks that reflect evolving societal values.
\end{itemize}

\section{Mathematical Framework}

\subsection{Modeling Social Structures}
Human networks can be represented as a tensor network \( \tensorNet{ij} \), where \( i \) and \( j \) denote individual nodes and their connections. The time-dependent state of a network is expressed as:
\begin{equation}
    \Psi(t) = \sum_{i,j} \tensorNet{ij}(t) \otimes \superpositionState{i}(t) \superpositionState{j}(t),
\end{equation}
where \( \superpositionState{i}(t) \) represents the superposition of states for node \( i \).

\subsection{Entanglement in Shared Decisions}
To model shared decisions, we define an entanglement term based on eigenvalues \( \eigenVal{i} \):
\begin{equation}
    E_{ij} = \eigenVal{i} \cdot \eigenVal{j} \cdot \cos(\theta_i - \theta_j),
\end{equation}
where \( \theta \) denotes the phase angle, capturing alignment in decision-making.

\subsection{Superposition in Outcomes}
Superposition is modeled using prime-based encoding \( \primeEncode{i} \) for nodes:
\begin{equation}
    \superpositionState{i} = \sum_{k} \frac{\primeEncode{k}}{Z} \ket{k},
\end{equation}
where \( Z \) is the normalization constant.

\subsection{Tunneling in Transitions}
Unexpected societal shifts can be captured by a non-linear multiplicity equation:
\begin{equation}
    \frac{\partial \rho_k}{\partial t} = \alpha_k \rho_k + \beta_k \sum_{j} T_{kj} \rho_j + \gamma_k \rho_k^2,
\end{equation}
where \( \rho_k \) represents the state density, \( T_{kj} \) models influence between nodes, and \( \alpha_k, \beta_k, \gamma_k \) are dynamic parameters.

\section{Applications and Experimentation}

\subsection{Quantifying Propagation}
Using eigenvalue analysis, the stability of influence propagation is examined through:
\begin{equation}
    \Delta \lambda_i = \feedbackFunc{i}(t) \cdot \lambda_i - \xi(t),
\end{equation}
where \( \feedbackFunc{i}(t) \) is the recursive feedback and \( \xi(t) \) models stochastic noise.

\subsection{Ethical Frameworks}
Dynamic ethical models integrate tensor networks and recursive feedback:
\begin{equation}
    \mathcal{E}(t) = \sum_{i} \tensorNet{ij}(t) \cdot \feedbackFunc{ij}(t),
\end{equation}
where \( \mathcal{E}(t) \) is the evolving ethical landscape.

\subsection{Simulating Societal Changes}
A real-time simulation incorporates stochastic noise and feedback loops:
\begin{equation}
    \Psi_{\text{society}}(t) = \sum_{i,j} \mathcal{T}_{ij}(t) \cdot e^{\mathrm{i}\phi(t)} \cdot (1 + \epsilon(t)),
\end{equation}
where \( \epsilon(t) \) represents random external influences.

\section{Conclusion}
By integrating Multiplicity Theory into Quantum Anthropology and Social Physics, we provide a framework to model complex societal dynamics and ethical evolution. Future work involves refining simulations and applying the models to large-scale societal datasets.


\nocite{*}
\bibliographystyle{plain}
\bibliography{references}

\end{document}